%% This file was a short chapter of its own until 2026-09-20.  It is now the closing part of the
%% chapter on the expression language, so that the chapter on proofs follows that chapter directly.
\section{Formulas as typed objects: \texttt{wff}}
\label{chap:syntax-contexts}

A \emph{well-formed formula} in this system is not a bare S-expression
but a typed object: a formula that has passed the well-formedness check.  In the REPL it prints as:
\begin{VNBsnip}
#[wff N "formula string"]
\end{VNBsnip}
where \texttt{N} is an internal identifier and the quoted string is the
pretty-printed formula.  Every formula in a sequent is a \texttt{wff}.

\texttt{(make-wff expr)} constructs such an object, from a string in
surface syntax or from an S-expression.  \texttt{(free-vars p)} returns
the variables that are free in it.

Two procedures expose the internals of a \texttt{wff}:
\begin{VNBsnip}
(wff->sexp p)        ; => raw S-expression (no output)
(display-contents p) ; prints kind, library name, formula
\end{VNBsnip}
\texttt{wff->sexp} is defined in \texttt{wff.scm};
\texttt{display-contents} is defined in \texttt{contexts.scm}.
Both are available in interactive sessions.

Example (the base library is \texttt{VNB-LIBRARY}):
\begin{VNBsnip}
(define p (make-wff "forall([x in A], 0 <= funny-name)"))
; => #[wff 12 "forall([x in a], 0 <= funny-name)"]

(free-vars p)
; => (a funny-name)   ; x is bound by the forall
\end{VNBsnip}
